\scanpage{658}\mspage{46}
Donc ici $u^W_!$ et $u^W_*$ n'existent pas, \unsure{pour}
$I\unsure{-|\to}e$ avec $I$ catégorie finie. Mais
aussi $M$ n'est pas stable par \inserted{petites} $\varinjlim$.

\uline{Exemple 3.12} Soit $M$ une catégorie munie
d'un localiseur $W$ (fortement saturé) tel
que $M-\underset{\inserted{\substack{\Vert\\\unsure{\mathscr H|\mathcal H}}}}{MW^{-1}}$ admette un adjoint
-- à gauche disons. On a donc une situation
de deux foncteurs adjoints
\diagram{\begin{minipage}{.92\linewidth}
\[
 \begin{tikzcd}[ampersand replacement=\&,column sep=6em]
 \mathscr H \arrow[r,shift left=1.0ex,"{v_!\ \text{pl. fid}}","{\scriptscriptstyle\vert}"{marking,pos=.55}] \&
 M \arrow[l,shift left=1.0ex,"{v^*\ \text{localisation}}"]
 \end{tikzcd}\,,\tag*{\unsure{(3.12 1)|(3.12.1)}}
\]
\end{minipage}}
et $W=(v^*)^{-1}(\mathrm{Iso}(\mathscr H))$. On trouve
que pour \unsure{tt|tte} $I$ dans Cat, \unsure{-- .|on a}
\[
 \mathrm H\deleted{\unsure{\mathrm O}}\mathrm o_W(I)\simeq\mathscr H(I)\overset{\mathrm{def}}{=}\mathrm{Hom}(I^\circ,\mathscr H).\tag*{\unsure{(3.17.2)|(3.12.2)}}
\]
On peut donc songer à faire un
exemple où dans $M$ les petites $\underset{\longleftarrow}{\varinjlim}$
existent, mais non dans $\mathscr H$ ; de sorte
que les $u^W_!$ ou $u^W_*$ n'existent pas
tjrs. Mais ça semble difficile à
construire, vu que dans la situation
envisagée, on voit facilement que les \inserted{petites}
$\varprojlim$ existent dans $\mathscr H$. Donc si dans $\mathscr H$
% archive sheet number: 657
